
\section{Effective inverse stability, adaptive exploration, and physical-time uncertainty}
\label{sec:quantitative-jacobi}

The response metric in Theorem~\ref{thm:jacobi-rate} is intrinsic, but a
coefficient theorem also requires a computable modulus.  We now construct one.
The bounds below are intentionally conservative: their role is to expose the
depth, separation, sampling-weight, and elapsed-time dependence rather than to
hide analytic continuation inside compactness.

Put
\[
 B_J=b_++2a_+,
 \qquad \Lambda_A=1+c_++B_J,
 \qquad Q=8\max\{2,\Lambda_A,T,T^{-1},a_-^{-1}\}.
\]
For \(J\ge0\), set \(R_J=4J+8\) and
\begin{equation}
 L_J=Q^{12(J+1)^3}.
 \label{eq:effective-LJ}
\end{equation}
All these constants are determined by the declared coefficient box and the
probe horizon.

\begin{lemma}[Effective response-jet inversion]
\label{lem:effective-response-jet}
For every \(J\ge0\) and \(\beta,\gamma\in\mathfrak B\),
\begin{equation}
 d_J(\beta,\gamma)
 \le L_J\max_{0\le r\le R_J}
       |h_\beta^{(r)}(0)-h_\gamma^{(r)}(0)|.
 \label{eq:effective-jet-inverse}
\end{equation}
The exponent in \eqref{eq:effective-LJ} is not optimized.
\end{lemma}

\begin{proof}
Let \(\mu_k(\beta)=\langle e_0,J_\beta^ke_0\rangle\).  Repeatedly using
\(A_\beta(q,v)=(v,-J_\beta q-cv)\) shows triangularly that
\(h_\beta^{(2k+2)}(0)=(-1)^k\mu_k(\beta)\) plus a polynomial in
\(c,\mu_0,\ldots,\mu_{k-1}\), while the odd derivative gives the companion
triangular relation.  Hence derivatives through order \(4J+8\) recover
\(c,\mu_0,\ldots,\mu_{2J+2}\) by additions and multiplications whose
Lipschitz constants are bounded by \(Q^{4(J+1)^2}\).

Write
\(\Delta_j=\det(\mu_{r+s})_{r,s=0}^j\).  The monic orthogonal-polynomial
identity gives
\begin{equation}
 \Delta_j=\prod_{r=0}^{j-1}a_r^{2(j-r)}
 \ge a_-^{j(j+1)}.
 \label{eq:hankel-determinant-lower}
\end{equation}
Moreover \(|\mu_k|\le B_J^k\).  The standard determinant formulas for the
Jacobi recurrence coefficients, followed by the adjugate formula and
Hadamard's inequality, therefore bound the Lipschitz constant of the map
\((\mu_0,\ldots,\mu_{2J+2})\mapsto(a_0,b_0,\ldots,a_J,b_J)\) by
\(Q^{8(J+1)^3}\).  Multiplying the two displayed bounds gives
\eqref{eq:effective-jet-inverse}.  This proof is an effective finite
calculation; in particular, \(L_J\) is not a minimum over a compact set.
\end{proof}

For \(R\ge1\), let
\[
 V_R=(k^r/r!)_{0\le k,r\le R},
 \qquad C_R=\|V_R^{-1}\|_\infty,
 \qquad M_R=2e^{\Lambda_AT}\Lambda_A^R.
\]
The number \(C_R\) is a completely explicit rational number.  Lagrange
interpolation at \(0,\ldots,R\) gives the useful closed bound
\begin{equation}
 C_R\le 2^{R+1}(R+1)!.
 \label{eq:vandermonde-inverse-bound}
\end{equation}
Indeed, the coefficient sum of
\(\prod_{0\le j\le R,\,j\ne k}(x-j)\) is at most
\((R+1)!/(k+1)\), and multiplication by
\(r!/[k!(R-k)!]\) bounds every entry of \(V_R^{-1}\); summing a row gives
\eqref{eq:vandermonde-inverse-bound}.  For
\(0<\varepsilon\le\min(1,T/R)\), define
\begin{equation}
 E_R(\varepsilon)=
 \frac{C_RM_RR^{R+1}}{(R+1)!}\,\varepsilon.
 \label{eq:taylor-grid-error}
\end{equation}

\begin{lemma}[Finite-grid derivative certificate]
\label{lem:finite-grid-certificate}
If \(g=h_\beta-h_\gamma\), then
\begin{equation}
 \max_{0\le r\le R}|g^{(r)}(0)|
 \le C_R\varepsilon^{-R}
       \max_{0\le k\le R}|g(k\varepsilon)|
       +E_R(\varepsilon).
 \label{eq:finite-grid-derivative}
\end{equation}
Here \(g(0)=0\) is known exactly, so only the positive grid points require
measurements.
\end{lemma}

\begin{proof}
For \(r\ge1\),
\(h_\beta^{(r)}(t)=\ell A_\beta^{r-1}e^{tA_\beta}B\); hence
\(\sup_{t\le T}|g^{(R+1)}(t)|\le
2e^{\Lambda_AT}\Lambda_A^R=M_R\).  Taylor's formula at the \(R+1\)
points \(k\varepsilon\) therefore gives
\(V_R(\varepsilon^rg^{(r)}(0))_{r=0}^R\) plus a remainder bounded by
\(M_R(R\varepsilon)^{R+1}/(R+1)!\).  Multiply by \(V_R^{-1}\) and use
\(\varepsilon^{-r}\le\varepsilon^{-R}\).
\end{proof}

We specialize the diagnostic law to a canonical union of two countable
families.  Half of its mass is any fixed explicit dense enumeration
\((u_m)\) of the rational points of \((0,T]\), with weight \(2^{-m-1}\).
The remaining half is assigned to the multiscale cells
\[
 q=(J,s,k),\qquad J\ge0,\ s\ge1,\ 1\le k\le R_J,
\]
with
\begin{equation}
 \varepsilon_{J,s}=\min\left\{1,\frac{T}{2R_J}\right\}2^{-s},\qquad
 t_q=k\varepsilon_{J,s},\qquad
 \omega_q=\frac{2^{-J-s-2}}{R_J}.
 \label{eq:multiscale-design}
\end{equation}
The weights sum to one and the dense half preserves the hypotheses of
Theorem~\ref{thm:jacobi-rate}.

For \(0<\delta\le1\), let \(s_J(\delta)\) be the least positive integer for
which
\begin{equation}
 E_{R_J}(\varepsilon_{J,s_J(\delta)})
 \le \frac{\delta}{2L_J},
 \label{eq:scale-choice}
\end{equation}
and put
\begin{align}
 \underline\kappa_J(\delta)
 &=\frac{2^{-J-s_J(\delta)-2}}{R_J}
 \left\{
   \frac{\delta\,\varepsilon_{J,s_J(\delta)}^{R_J}}
        {2C_{R_J}L_J}
 \right\}^2,
 \label{eq:effective-kappa}\\
 \zeta_R&=\min\{1,T/(2R)\},\qquad
 H_R=\frac{C_RM_RR^{R+1}}{(R+1)!},\notag\\
 G_J&=\max\{2/\zeta_{R_J},\,4L_JH_{R_J}\},\notag\\
 K_J&=\frac{2^{-J-4}}
 {R_J\zeta_{R_J}C_{R_J}^2L_J^2G_J^{2R_J+1}}.
 \label{eq:explicit-kappa-prefactor}
\end{align}

\begin{theorem}[Explicit depth-dependent Jacobi stability]
\label{thm:effective-jacobi-stability}
For every \(J\ge0\), \(0<\delta\le1\), and
\(\beta,\gamma\in\mathfrak B\),
\begin{equation}
 d_J(\beta,\gamma)\ge\delta
 \quad\Longrightarrow\quad
 D(\beta,\gamma)\ge\underline\kappa_J(\delta)>0.
 \label{eq:effective-response-separation}
\end{equation}
Consequently \(\kappa_J(\delta)\ge\underline\kappa_J(\delta)\), and the
fully explicit lower estimate
\begin{equation}
 \underline\kappa_J(\delta)\ge K_J\delta^{2R_J+3}
 \ge \exp\{-A_0(J+1)^4\}\delta^{20(J+1)},
 \qquad 0<\delta\le1,
 \label{eq:kappa-asymptotic-lower}
\end{equation}
holds with the displayed prefactor \(K_J\) and, for example,
\begin{equation}
 A_0=10^6\{1+\Lambda_AT+|\log T|+\log Q\}.
 \label{eq:explicit-A0}
\end{equation}
\end{theorem}

\begin{proof}
If \(d_J\ge\delta\), Lemma~\ref{lem:effective-response-jet} gives a
derivative difference at least \(\delta/L_J\).  Apply
Lemma~\ref{lem:finite-grid-certificate} at the scale selected in
\eqref{eq:scale-choice}.  Some measured grid value then has magnitude at
least
\(\delta\varepsilon_{J,s}^{R_J}/(2C_{R_J}L_J)\).  Its single contribution
to \(D\) is exactly the right side of \eqref{eq:effective-kappa}.
To prove the first inequality in \eqref{eq:kappa-asymptotic-lower}, write
\(\varepsilon_{J,s}=\zeta_{R_J}2^{-s}\).  If the least admissible scale has
\(s=1\), then \(\varepsilon_{J,s}=\zeta_{R_J}/2\ge\delta/G_J\).  If
\(s>1\), minimality says
\(2H_{R_J}\varepsilon_{J,s}>\delta/(2L_J)\), and again
\(\varepsilon_{J,s}>\delta/G_J\).  Hence
\(2^{-s}=\varepsilon_{J,s}/\zeta_{R_J}
\ge\delta/(\zeta_{R_J}G_J)\).  Substitution into
\eqref{eq:effective-kappa} gives exactly
\(\underline\kappa_J(\delta)\ge K_J\delta^{2R_J+3}\).
Finally \eqref{eq:vandermonde-inverse-bound},
\(R_J\le8(J+1)\), \(\log x\le x\), and the definitions of
\(L_J,H_R,G_J\) give
\(-\log K_J\le A_0(J+1)^4\) with \(A_0\) in
\eqref{eq:explicit-A0}; also \(2R_J+3=8J+19\le20(J+1)\).
Since \(0<\delta\le1\), the second inequality follows.  Thus no constant in
the inverse modulus is defined by an unnamed extremum or compactness
argument.
\end{proof}

\begin{corollary}[Coefficient posterior exponent with no hidden modulus]
\label{cor:effective-cylinder-rate}
Under the multiscale design,
\begin{equation}
 \limsup_{n\to\infty}\frac1n
 \log\Pi_n^{\nu_n}\{d_J(\beta,\beta_0)\ge\delta\}
 \le-\frac{a^2}{2\sigma^2}\underline\kappa_J(\delta)
 \quad\text{almost surely}.
 \label{eq:effective-cylinder-posterior}
\end{equation}
\end{corollary}

\begin{proof}
Combine Theorem~\ref{thm:jacobi-rate} with
Theorem~\ref{thm:effective-jacobi-stability}.
\end{proof}

\subsection{Finite propagation and increasing-depth recovery}

\begin{lemma}[Boundary locality for a tridiagonal bath]
\label{lem:jacobi-boundary-locality}
There is an explicit \(C_{\rm loc}\), depending only on the coefficient box,
such that if \(\beta\) and \(\gamma\) agree through depth \(K\), including
\(c\), then
\begin{equation}
 \|h_\beta-h_\gamma\|_{C([0,T])}
 \le C_{\rm loc}e^{\Lambda_AT}
       \frac{(\Lambda_AT)^{K+1}}{(K+1)!}.
 \label{eq:factorial-boundary-locality}
\end{equation}
If their first \(K\) coefficients differ by at most \(r\), the right side is
augmented by \(C_{\rm loc}e^{\Lambda_AT}(K+1)r\).
\end{lemma}

\begin{proof}
Expand the two bounded semigroups in powers of their generators.  A word that
starts at the boundary, visits a coefficient below depth \(K\), and returns
to the boundary contains at least \(K+1\) nearest-neighbour transmissions.
All shorter words cancel.  Summing the remaining exponential-series tail
gives \eqref{eq:factorial-boundary-locality}.  Telescoping the finitely many
changed boundary coefficients gives the Lipschitz term.
\end{proof}

Assume for the next result that the Jacobi prior is a product prior whose
one-coordinate densities are bounded below by one common \(p_->0\) on their
compact intervals.  Let \(\mathcal H(\epsilon)\) denote the logarithm of the
\(C([0,T])\) covering number of the response image at radius \(\epsilon\).
Let \(K(\epsilon)\) be the least integer for which the factorial term in
\eqref{eq:factorial-boundary-locality} is at most \(\epsilon/4\), put
\[
 A_K=C_{\rm loc}e^{\Lambda_AT}(K+1),
 \qquad
 \mathcal P(\epsilon)=(2K(\epsilon)+3)
   \log\frac{4A_{K(\epsilon)}}{p_-\epsilon},
\]
and interpret \(\mathcal P\) as a deterministic lower bound on minus the log
prior mass of an \(\epsilon\)-response ball.  Lemma~\ref{lem:jacobi-boundary-locality}
gives the effective bounds
\begin{equation}
 \mathcal H(\epsilon)+\mathcal P(\epsilon)
 \le C\frac{|\log\epsilon|^2}{\log(2+|\log\epsilon|)}
 \qquad(0<\epsilon<1/2).
 \label{eq:response-entropy-bound}
\end{equation}

\begin{theorem}[Growing-depth posterior recovery]
\label{thm:growing-depth-recovery}
Let \(J_n\uparrow\infty\), \(0<\delta_n\le1\), and put
\(k_n=\underline\kappa_{J_n}(\delta_n)\).  If
\begin{equation}
 n k_n^2-\mathcal H(k_n/16)\longrightarrow+\infty,
 \qquad
 \mathcal P(\sqrt{k_n}/8)=o(nk_n),
 \label{eq:growing-depth-conditions}
\end{equation}
then
\begin{equation}
 \Pi_n^{\nu_n}\{d_{J_n}(\beta,\beta_0)\ge\delta_n\}
 \longrightarrow0
 \quad\text{in }P_{\beta_0,z_0}\text{-probability}.
 \label{eq:growing-depth-contraction}
\end{equation}
For every fixed \(\delta>0\), the concrete choice
\(J_n=o((\log n)^{1/5})\) satisfies the conditions.
\end{theorem}

\begin{proof}
Use an \(k_n/16\)-net of the compact response image.  Gaussian concentration
at the net points and sup-norm interpolation give a uniform likelihood error
smaller than \(k_n/8\) with probability tending to one under the first
condition.  The product-prior lower bound, combined with
Lemma~\ref{lem:jacobi-boundary-locality}, supplies a response ball of radius
\(\sqrt{k_n}/8\) whose negative log mass is \(o(nk_n)\) under the second
condition.  On the
complement of the cylinder, Theorem~\ref{thm:effective-jacobi-stability}
gives response information at least \(k_n\), so numerator over denominator
tends to zero.  Finally \eqref{eq:kappa-asymptotic-lower} and
\eqref{eq:response-entropy-bound} verify the stated concrete regime.
\end{proof}

For \(\tau>0\), let \(W_\tau e_j=e^{-\tau j}e_j\).

\begin{corollary}[Weighted operator-norm recovery]
\label{cor:weighted-operator-recovery}
There is a box constant \(C_{\rm box}\) such that
\begin{equation}
 \|W_\tau(J_\beta-J_{\beta_0})W_\tau\|_{\rm op}
 \le3d_J(\beta,\beta_0)+C_{\rm box}e^{-2\tau(J+1)}.
 \label{eq:weighted-operator-bound}
\end{equation}
Hence any sequences in Theorem~\ref{thm:growing-depth-recovery} with
\(\delta_n\to0\) yield posterior contraction in this weighted operator norm.
\end{corollary}

\begin{proof}
Apply the Schur row-sum bound to the tridiagonal difference.  The first
\(J+1\) rows contribute at most \(3d_J\); every remaining matrix entry is
multiplied by at most \(e^{-2\tau(J+1)}\), and the coefficient box bounds its
row sum by \(C_{\rm box}\).
\end{proof}

\subsection{Adaptive diagnostic allocation}

The fixed multiscale law is not essential.  At stage \(i\), first draw a fresh
Bernoulli exploration flag \(E_i\) of parameter \(\rho_J\in(0,1]\).  If
\(E_i=1\), draw the diagnostic cell from \(\omega\); otherwise an arbitrary
predictable reward-seeking or stabilizing controller may choose the diagnostic
cell, in addition to choosing the preceding bounded feedback segment.  The
fresh sign and readout noise are still drawn after that choice.

\begin{theorem}[Exploration-floor adaptive identification]
\label{thm:adaptive-exploration-floor}
For every closed \(F\subset\mathfrak B\), the exact posterior under the
exploration-floor policy satisfies
\begin{equation}
 \limsup_{n\to\infty}\frac1n\log\Pi_n^{\nu_n}(F)
 \le-\frac{\rho_Ja^2}{2\sigma^2}
       \inf_{\beta\in F}D(\beta,\beta_0)
 \quad\text{almost surely}.
 \label{eq:adaptive-exploration-rate}
\end{equation}
Thus all effective coefficient and weighted-operator conclusions remain valid
with the exponent multiplied by \(\rho_J\), while the nonexploration actions
may be fully adaptive.
\end{theorem}

\begin{proof}
Conditionally on the past, the exploration contribution to every squared
response difference is \(\rho_JD(\beta,\beta_0)\); all other diagnostic
contributions are nonnegative.  The signed score is a bounded-envelope
Gaussian martingale field.  The same finite sup-norm nets used in
Lemma~\ref{lem:supnorm-response-slln} give a uniform martingale strong law.
The denominator has exponential rate zero by full support and uniform response
continuity.  The standard numerator--denominator argument then gives
\eqref{eq:adaptive-exploration-rate}.
\end{proof}

\subsection{Honest coefficient cylinders and elapsed-time rate}

Fix a finite set \(\mathcal Q\) of multiscale cells.  Let \(N_q(n)\) be the
number of visits to cell \(q\), and, when \(N_q(n)>0\), define
\[
 \widehat h_q=\frac1{aN_q(n)}
   \sum_{i\le n:M_i=q}S_iY_i.
\]
Let \(r_i=C(i+1)^{-1-\epsilon_w}\) be the deterministic envelope in
\eqref{eq:polynomial-washout-envelope}, and put
\begin{equation}
 R_{q,n}(\alpha)=
 \frac{\sigma}{a}\sqrt{\frac{2\log(2n|\mathcal Q|/\alpha)}{N_q(n)}}
 +\frac1{aN_q(n)}\sum_{i\le n:M_i=q}r_i,
 \label{eq:response-confidence-radius}
\end{equation}
with \(R_{q,n}=+\infty\) if \(N_q(n)=0\).

\begin{theorem}[Honest growing finite-block uncertainty]
\label{thm:honest-jacobi-cylinders}
The response confidence set
\begin{equation}
 \mathcal C_n(\alpha)=\{\beta\in\mathfrak B:
 |h_\beta(t_q)-\widehat h_q|\le R_{q,n}(\alpha)
 \text{ for every }q\in\mathcal Q\}
 \label{eq:response-confidence-set}
\end{equation}
has frequentist coverage at least \(1-\alpha\), uniformly over the coefficient
box, bounded feedback segments, and working initial-state priors.  If
\(\mathcal Q\) contains the grid selected by \(s_J(\delta)\) and
\(2\max_{q\in\mathcal Q}R_{q,n}(\alpha)\) is below the grid threshold in the
proof of Theorem~\ref{thm:effective-jacobi-stability}, then
\(\operatorname{diam}_{d_J}\mathcal C_n(\alpha)\le\delta\).  The statement
continues to hold for the growing regimes of
Theorem~\ref{thm:growing-depth-recovery} whenever the corresponding cell
counts diverge.
\end{theorem}

\begin{proof}
For every cell, the visit indicator at stage \(i\) is predictable before the
fresh Gaussian readout.  Therefore the selected Gaussian sums are martingale
transforms.  At each possible visit count \(1\le m\le n\), the corresponding
sum is \(N(0,m\sigma^2)\); a union bound over cells, signs, and the at most
\(n\) possible counts gives the first term of
\eqref{eq:response-confidence-radius}.  This argument remains valid when
future visits depend on earlier readouts.  The persistent-state bias is
bounded by the second term.  Apply
Lemma~\ref{lem:finite-grid-certificate} to any two elements of the confidence
set and then Lemma~\ref{lem:effective-response-jet} to obtain the diameter
claim.
\end{proof}

Finally, let \(\mathsf T_n\) be total elapsed physical time through stage
\(n\).  Bounded feedback and probe durations, together with
\eqref{eq:logarithmic-washout}, give
\begin{equation}
 \mathsf T_n=\kappa_w n\log n+O(n),
 \qquad \kappa_w=\frac{1+\epsilon_w}{\lambda_*}.
 \label{eq:physical-time-complexity}
\end{equation}
Let \(W\) denote the principal Lambert function and define
\begin{equation}
 v(t)=\frac{t}{\kappa_w W(t/\kappa_w)}
 \sim\frac{t}{\kappa_w\log t}.
 \label{eq:physical-time-speed}
\end{equation}
Then \(v(\mathsf T_n)/n\to1\).  Consequently every upper and lower posterior
large-deviation bound in Theorem~\ref{thm:jacobi-rate} remains valid in elapsed
physical time after replacing the stage speed \(n\) by
\(v(\mathsf T_n)\).  In particular the experiment has near-linear, rather
than quadratic, physical-time cost.
